\chapter{Frames and Notation}

\section{Frames}

The first implementation uses three frames:

\begin{center}
\begin{tabular}{lll}
\toprule
Symbol & Frame & Meaning \\
\midrule
$i$ & inertial & non-rotating inertial reference frame \\
$e$ & ECEF/PCPF & Earth-fixed frame rotating with Earth \\
$b$ & body/IMU & vehicle body frame, collocated with the IMU and navigation state \\
\bottomrule
\end{tabular}
\end{center}

The body frame $b$ is the first implementation's navigation state point. This
collocation assumption intentionally removes the origin-frame and IMU-lever-arm
terms derived in the broader navigation reference. A non-IMU aiding sensor
$s$, such as a GNSS antenna, may still have a configured body-frame lever arm
$l_{bs}^{b}$ from the body/IMU point to the sensor measurement point
\cite{groves2013}.

\section{Vector and Transform Conventions}

The nominal navigation variables are:
\begin{align}
\pb_{eb}^{e} &\in \R^3 && \text{position of body with respect to ECEF, resolved in ECEF},\\
\vb_{eb}^{e} &\in \R^3 && \text{velocity of body with respect to ECEF, resolved in ECEF},\\
\quat_b^e &\in \R^4 && \text{unit quaternion mapping body components into ECEF components}.
\end{align}

Let $C_b^e(\quat_b^e)$ denote the DCM corresponding to $\quat_b^e$ that maps
body-resolved components into ECEF-resolved components:
\begin{equation}
\mathbf{a}^{e}=C_b^e\mathbf{a}^{b}.
\end{equation}
The inverse mapping is
\begin{equation}
C_e^b=\left(C_b^e\right)^{T}.
\end{equation}

\begin{warningbox}
Code identifiers sometimes use explicit \texttt{from2to} names such as
\texttt{q\_e2b} or \texttt{rpy\_e2b\_rad} when transform direction is the
primary meaning. This document's mathematical notation follows the compact
Groves-style convention: $C_b^e$ maps from $b$ to $e$.
\end{warningbox}

\section{Truth, Estimate, Measurement, and Error Notation}

Unadorned variables represent the true physical quantity. A hat denotes the
current nominal estimate, and a tilde denotes a measured sensor quantity:
\begin{align}
\pb_{eb}^{e} & && \text{true position},&
\hat{\pb}_{eb}^{e} & && \text{estimated/nominal position},\\
\Delta\tilde{\boldsymbol\theta}_{ib}^{b} & && \text{measured gyro angle increment},&
\Delta\tilde{\mathbf{v}}_{ib}^{b} & && \text{measured accelerometer velocity increment}.
\end{align}

The additive error-state convention is truth minus estimate:
\begin{align}
\delta\pb^{e} &= \pb_{eb}^{e}-\hat{\pb}_{eb}^{e},\\
\delta\vb^{e} &= \vb_{eb}^{e}-\hat{\vb}_{eb}^{e},\\
\delta\bb_g &= \bb_g-\hat{\bb}_g,\\
\delta\bb_a &= \bb_a-\hat{\bb}_a.
\end{align}

For attitude, the error state is a three-component small rotation vector
resolved in ECEF. Following the convention used in the broader NavKit
reference, the true body-to-ECEF DCM is related to the nominal DCM by
\begin{equation}
C_b^{e,true}
\approx
\left(\I-\crossmat{\delta\boldsymbol\theta}\right)
\hat{C}_b^e.
\label{eq:v1_attitude_error_dcm_relation}
\end{equation}
This convention is intentionally selected to match the left-multiplicative
small-quaternion correction in \cref{chap:filter_correction_reset}
\cite{groves2013,sola2018}.

\section{Earth Rate}

The ECEF frame rotates with respect to the inertial frame at
\begin{equation}
\omegab_{ie}^{e}
=
\begin{bmatrix}
0 & 0 & \omega_{ie}
\end{bmatrix}^{T}.
\end{equation}
For Earth/WGS-84, $\omega_{ie}$ is supplied by the selected planet policy. The
Navigator must not know that value directly; the concrete mechanization or
propagation policy owns the planet and gravity configuration.
